\section{Evidence constructors and positive introspection in \(\PA\)}
\label{sec:p05-pa-positive-introspection}

We now construct the proof of the original evidence-bearing sentence. The
construction uses the selected formula
\(\mathsf{Prf}^{19}_{\PA}(p,x)\), namely the literal
\(\Gamma_{d_{19}}(p,x,\bar1)\), throughout.
The bounded universal constructor of
\cref{thm:p05-bounded-case-elimination} accepts an arbitrary formula as
its matrix. This permits a bounded range of existentially witnessed trace
steps to become one derivation of the original quantified trace condition.
The complete verifications are in
\cref{sec:annex-p05-pa-positive-introspection}.

\begin{hbdefinition}{def:p05-positive-evidence-formation}{Literal positive formation and its evidence}
A \emph{positive formation} here is a fixed marked arithmetic formula
obtained from bounded formulas by conjunction, disjunction, unbounded
existential quantification, and either marked bounded quantifier. The
bounded leaves may themselves contain any of the bounded-formula
connectives. At every node take the bounded-leaf clause first whenever
that entire subformula is bounded. This is an external parsing convention
for a fixed formula, with no numerical input ranging over formulas.

For such a formula \(\theta\), define its bounded evidence formula
\(\mathcal E_\theta(\bar x,z)\) by the following clauses. The tuple
\(\bar x\) lists the free parameters apart from a displayed binder.
The evidence variable and all auxiliary bound variables are chosen by
the fresh-variable convention of \cref{thm:p05-pa-sigma-evidence}.
Original formula binders and comparison binders are retained.
\[
\begin{aligned}
 \mathcal E_\theta(\bar x,z)
   &\defeq z=0\land\theta(\bar x)
                         &&\text{if \(\theta\) is bounded},\\
 \mathcal E_{\psi\land\chi}(\bar x,z)
   &\defeq\exists u\leq z\,\exists v\leq z\,
       (\mathsf P(u,v,z)\land\mathcal E_\psi(\bar x,u)
                              \land\mathcal E_\chi(\bar x,v)),\\
 \mathcal E_{\psi\lor\chi}(\bar x,z)
   &\defeq\exists a<\bar2\,\exists v\leq z\,
       [\mathsf P(a,v,z)\land
        ((a=0\land\mathcal E_\psi(\bar x,v))
          \lor(a=\bar1\land\mathcal E_\chi(\bar x,v)))],\\
 \mathcal E_{\exists x_i\psi}(\bar x,z)
   &\defeq\exists x_i\leq z\,\exists v\leq z\,
        (\mathsf P(x_i,v,z)\land\mathcal E_\psi(\bar x,x_i,v)).
\end{aligned}
\]
For \(\diamond\in\{<,\leq\}\), the two marked clauses are
\[
\begin{aligned}
 \mathcal E_{\exists x_i\mathrel\diamond t\,\psi}(\bar x,z)
  &\defeq\exists x_i\mathrel\diamond t(\bar x)\,
       \exists v\leq z\,
       (\mathsf P(x_i,v,z)\land\mathcal E_\psi(\bar x,x_i,v)),\\
 \mathcal E_{\forall x_i\mathrel\diamond t\,\psi}(\bar x,z)
  &\defeq\forall x_i\mathrel\diamond t(\bar x)\,
       \exists v\leq z\,
       (\mathsf B(z,x_i,v)\land\mathcal E_\psi(\bar x,x_i,v)).
\end{aligned}
\]
The bounding term has no occurrence of its bound variable. The displayed
\(\mathsf P\) and \(\mathsf B\) are the original formulas of
\cref{lem:pairing-graph,lem:goedel-beta}. These clauses do not convert
the original formula to another syntax tree. When it is already in the
fixed marked negation normal form, \(\mathcal E_\theta\) is exactly
the \(\mathsf{Ev}_\theta\) of
\cref{thm:p05-pa-sigma-evidence}, with the same auxiliary choices.

List each fixed formula's distinct free indices in increasing order.
Numerical inputs \(\bar m\) replace them by \(\nu(m_j)\), in that
order, using the selected substitution description. Denote the resulting
code by \(\theta^\nu(\bar m)\). A child receives its parent tuple
by the fixed free-index projections; a quantified child additionally
receives the current numerical witness at index \(i\). An inner binder
of the same name overrides that entry only within its own scope.
Write
\[
       e_\theta(\bar m,z)\defeq
                  b_{\mathcal E_\theta}(\bar m,z)
\]
for the zero-one description already specified for this fixed bounded
evidence formula by \cref{def:p05-fixed-bounded-builders}. In this
application, fixed tuple projections place the original parameter values
and the final evidence value in the increasing free-index order of
\(\mathcal E_\theta\); the fresh evidence index need not be last in
that order. The displayed argument order of \(e_\theta\) itself is
the original parameter tuple followed by its evidence value. Function
applications always expand through their specified \(\Gamma\) graphs.
\end{hbdefinition}

\begin{hbdefinition}{def:p05-evidence-projections}{Pair and beta evidence projections}
Using the already fixed head and tail descriptions, put
\[
       \pi_0(z)\defeq h(Sz),\qquad \pi_1(z)\defeq t(Sz).
\]
These invert \(\mathsf P\), whose numerical pairing is one less than
the positive cons code. They are not head and tail of \(z\) itself.
Let \(b_{\mathsf B}(z,k,v)\) be the bounded Boolean description of the
literal fixed beta formula, with its free inputs in the displayed order
by the appropriate tuple projections. Define
\[
 \beta_*(z,k)\defeq
       \text{the least }v<Sz\text{ with }b_{\mathsf B}(z,k,v)=1,
       \quad\text{with failure value }0.
\]
This is exactly the bounded least-search fold of
\cref{def:p05-checker-description-language}: start at \([0,0]\),
record the first successful counter, and return its second field after
\(Sz\) updates. The bound is fixed before the fold.
\end{hbdefinition}

\begin{hblemma}{thm:p05-evidence-projections}{proved=annex:p5.evidence-projections}{Internal evidence decoding and existence}
\(\PA\) proves, with universal numerical closure,
\[
\begin{aligned}
 &\mathsf P(\pi_0(z),\pi_1(z),z),\qquad
                  \pi_0(z),\pi_1(z)\leq z,\\
 &\mathsf P(u,v,z)\to(u=\pi_0(z)\land v=\pi_1(z)),\\
 &\mathsf B(z,k,\beta_*(z,k)),\qquad \beta_*(z,k)\leq z,\\
 &\mathsf B(z,k,v)\to v=\beta_*(z,k).
\end{aligned}
\]
For every fixed positive formation \(\theta\), its evidence formula
is bounded and \(\PA\) proves
\[
 \forall\bar x\,
       (\theta(\bar x)\leftrightarrow
                         \exists z\,\mathcal E_\theta(\bar x,z)),
 \qquad
 \forall\bar x\forall z\,
       (e_\theta(\bar x,z)=1\leftrightarrow
                              \mathcal E_\theta(\bar x,z)).
\]
These assertions include zero evidence, an empty parameter tuple, empty
strict ranges, weak ranges with endpoint zero, and dependent inner witnesses.
\end{hblemma}

\begin{hbdefinition}{def:p05-positive-evidence-builders}{The exact proof description at each evidence clause}
Use the raw templates of
\cref{def:p05-builder-row-templates,def:p05-bounded-builder-conventions}
and the safe-tree invariant of
\cref{def:p05-bounded-builder-conventions}. We specify one raw description
\(V_\theta(\bar m,z)\) for each fixed positive formation. Every clause
below has the outer guard \(e_\theta(\bar m,z)=1\); its value is zero
when the guard fails. On success it has the following value. The letters
\(c,d\) in a connective clause denote the respective child numeral codes.

For a bounded leaf use
\[
                  V_\theta(\bar m,z)=R_\theta^+(\bar m).
\]
For conjunction and disjunction put \(a=\pi_0(z),v=\pi_1(z)\), and use
\[
\begin{aligned}
 V_{\psi\land\chi}(\bar m,z)
   &=\mathsf{And}(V_\psi(\bar m,a),V_\chi(\bar m,v)),\\
 V_{\psi\lor\chi}(\bar m,z)
   &=\begin{cases}
       \mathsf{OL}(V_\psi(\bar m,v),d),&a=0,\\
       \mathsf{OR}(c,V_\chi(\bar m,v)),&a\ne0.
     \end{cases}
\end{aligned}
\]
Child tuple projections are the fixed projections declared above; the
notation \(\bar m\) in each child suppresses only those projections.
The second disjunction branch is used under the outer success guard,
which will imply \(a=1\).

For \(\theta=\exists x_i\psi\), put \(a=\pi_0(z),v=\pi_1(z)\).
Let \(c\) be the literal matrix code after numeral substitution for its
other free indices, leaving index \(i\) available. Use
\[
 V_\theta(\bar m,z)=
          \mathsf{EI}(V_\psi(\bar m,a,v),i,c,\nu(a)).
\]
In \(V_\psi(\bar m,a,v)\), the tuple consists of the original
parameters together with the witness \(a\); its final argument is the
evidence value \(v\). An unused index receives no tuple position.

For a marked existential or universal, put
\[
 u=t^\nu(\bar m),\quad b=t^\#(\bar m),\quad
 N=N_\diamond(b),\quad G=G_\diamond(i,h,u),
\]
where \(h\) is the comparison binder of that original marked occurrence.
The literal matrix code \(c\) again leaves only the possible free index
\(i\). The safe bound-equality tree is \(R_t(\bar m)\).
For each numerical \(a\) in the range, let \(L(\bar m,a)\) denote
the raw positive literal tree for
\(\mathsf{sub}(G,i,\nu(a))\). It is the raw record of the already
specified fixed literal builder, decoded by the exact outer tags 19 and
15, with failure value zero. Its supplied comparison binder is \(h\).

For \(\theta=\exists x_i\mathrel\diamond t\,\psi\), set
\(a=\pi_0(z),v=\pi_1(z)\), and return
\[
 \mathsf{EI}\bigl(
   \mathsf{And}(L(\bar m,a),V_\psi(\bar m,a,v)),
                                    i,G\mathbin{\&}c,\nu(a)\bigr).
\]
For \(\theta=\forall x_i\mathrel\diamond t\,\psi\), define the
list description, with counter last,
\[
\begin{aligned}
 Q_\theta(\bar m,z,0)&=0,\\
 Q_\theta(\bar m,z,Sk)&=
    \mathsf{cat}(Q_\theta(\bar m,z,k),
                  [V_\psi(\bar m,k,\beta_*(z,k))]).
\end{aligned}
\]
Return the actual quantified tree
\[
 V_\theta(\bar m,z)=
       \mathsf U_\diamond
          (c,i,h,u,b,R_t(\bar m),Q_\theta(\bar m,z,N)).
\]
This is the \(\mathsf U_\diamond\) of
\cref{def:p05-bounded-case-elimination}; the matrix \(c\) need not
be bounded. Its finite-case equality transports, range-exhaustion proof,
original-guard discharge and universal introduction are part of this
expression. The list alone is not its output proof.

Finally let \(d_\theta\) be the finite description
\[
 d_\theta(\bar m,z)\defeq
  I(e_\theta(\bar m,z),\mathsf{pack}(V_\theta(\bar m,z)),0).
\]
The same symbol denotes this specified formation tree when used as the
subscript of \(\Gamma\). The extra outer guard preserves the zero
failure value after packing. Clauses are expanded through the previous
library in their displayed order; each numerical recursion calls only
already fixed child descriptions and its preceding state. Parentheses
and the original formula tree fix every conjunction association.
\end{hbdefinition}

\begin{hbtheorem}{thm:p05-positive-evidence-builders}{proved=annex:p5.positive-evidence-builders}{Evidence produces the literal sentence proof inside PA}
For every fixed positive formation \(\theta\), \(\PA\) proves
\[
 \forall\bar m\forall z\,
  (\mathcal E_\theta(\bar m,z)\to
           \mathsf{Safe}(V_\theta(\bar m,z),\theta^\nu(\bar m))).
\]
It proves the sentence test for every such numeral-instance root, and hence
\[
 \forall\bar m\forall z\,
 \bigl(\mathcal E_\theta(\bar m,z)\to
   \exists p\,(\Gamma_{d_\theta}(\bar m,z,p)\land
                   \mathsf{Prf}^{19}_{\PA}(p,\theta^\nu(\bar m)))\bigr).
\]
For a fixed closed positive formation \(\sigma\), the parameter tuple
is empty and its root is the fixed original code \(\ulcorner\sigma\urcorner\).
Consequently
\[
 \PA\derives\sigma\to
             \mathsf{Prv}^{19}_{\PA}
                         (\overline{\ulcorner\sigma\urcorner}).
\]
All these statements concern the same literal selected checker graph.
\end{hbtheorem}

\begin{hbdefinition}{def:p05-signed-evidence-builders}{Original syntax from signed normal-form evidence}
For a fixed formula \(\theta\), write \(\mathcal N^+\theta\)
for its fixed marked negation normal form and \(\mathcal N^-\theta\)
for the corresponding conversion of \(\neg\theta\). They have the
following literal formation. At an equality or bottom use \(\theta\)
and \(\neg\theta\); at a negation exchange the two child forms.
At the three binary connectives use
\[
\begin{array}{c|c|c}
 \theta&\mathcal N^+\theta&\mathcal N^-\theta\\\hline
 \psi\land\chi&\mathcal N^+\psi\land\mathcal N^+\chi
                         &\mathcal N^-\psi\lor\mathcal N^-\chi\\
 \psi\lor\chi&\mathcal N^+\psi\lor\mathcal N^+\chi
                         &\mathcal N^-\psi\land\mathcal N^-\chi\\
 \psi\to\chi&\mathcal N^-\psi\lor\mathcal N^+\chi
                         &\mathcal N^+\psi\land\mathcal N^-\chi.
\end{array}
\]
At a quantifier the plus form retains the quantifier and uses the plus
child; the minus form uses its dual and the minus child. A marked
quantifier retains its original bound, binder and comparison binder
in both forms. These rewrites move no formula across a binder, so
they create no capture requiring a new bound name.

Consider only a separately fixed pair \((\theta,\epsilon)\),
\(\epsilon\in\{+,-\}\), for which
\(\mathcal N^\epsilon\theta\) is a \(\Sigma^0_1\) positive
formation. The sign is a metalinguistic choice, not a numerical input.
Set
\[
 \mathcal E_\theta^\epsilon\defeq
       \mathcal E_{\mathcal N^\epsilon\theta},\qquad
 e_\theta^\epsilon\defeq b_{\mathcal E_\theta^\epsilon}.
\]
Use exactly the evidence occurrence and auxiliary choices produced by
the original fixed normal-form evidence construction. Equal-looking
subformulas at different occurrences retain their respective choices.
Inputs are the original free-parameter tuple followed by evidence;
the Boolean call uses its fixed free-index projections as before.

Define \(\widehat V_\theta^\epsilon(\bar m,z)\) as follows.
Every clause has outer success guard \(e_\theta^\epsilon(\bar m,z)=1\)
and value zero when it fails. Take the bounded-original-formula clause
first. On that clause use \(R_\theta^+\) for sign plus and
\(R_\theta^-\) for sign minus. In all remaining clauses put
\(a=\pi_0(z),v=\pi_1(z)\), and let \(c,d\) be the original
child numeral codes. A child call below always uses its own required
sign and its fixed parameter projections.

At \(\theta=\neg\psi\), the plus tree is
\(\widehat V_\psi^-(\bar m,z)\), and the minus tree is
\[
 \mathsf{NI}\bigl(
   \mathsf B(\widehat V_\psi^+(\bar m,z),
                         \mathsf L(\mathsf{Neg}(c),1)),
                                      \mathsf{Neg}(c),0\bigr).
\]
For positive conjunction use
\(\mathsf{And}(\widehat V_\psi^+(\bar m,a),
\widehat V_\chi^+(\bar m,v))\). For negative conjunction put
\(H=\mathsf L(c\mathbin{\&}d,1)\), and use
\[
 \begin{cases}
  \mathsf{NI}(\mathsf B(\mathsf P_0(H),
        \widehat V_\psi^-(\bar m,v)),c\mathbin{\&}d,0),&a=0,\\
  \mathsf{NI}(\mathsf B(\mathsf P_1(H),
        \widehat V_\chi^-(\bar m,v)),c\mathbin{\&}d,0),&a\ne0.
 \end{cases}
\]
Here the evidence pair is a disjunction tag and payload. For positive
disjunction use
\(\mathsf{OL}(\widehat V_\psi^+(\bar m,v),d)\) if \(a=0\),
and \(\mathsf{OR}(c,\widehat V_\chi^+(\bar m,v))\) otherwise.
For negative disjunction its evidence pair instead contains the two
child evidence values. Put
\[
 K_\psi=\mathsf B(\mathsf L(c,1),\widehat V_\psi^-(\bar m,a)),\qquad
 K_\chi=\mathsf B(\mathsf L(d,1),\widehat V_\chi^-(\bar m,v)),
\]
and return
\[
 \mathsf{NI}(\mathsf{OE}(\mathsf L(c\mathbin{|}d,1),
                            K_\psi,K_\chi,0,0),c\mathbin{|}d,0).
\]
For positive implication put \(H=\mathsf L(c,1)\), and use
\[
 \begin{cases}
  \mathsf{II}(\mathsf{EF}(\mathsf B(H,
                  \widehat V_\psi^-(\bar m,v)),d),c,0),&a=0,\\
  \mathsf{II}(\mathsf{Keep}(H,
                  \widehat V_\chi^+(\bar m,v)),c,0),&a\ne0.
 \end{cases}
\]
For negative implication return
\[
 \mathsf{NI}\bigl(
   \mathsf B(\mathsf{MP}^{\rm row}(\widehat V_\psi^+(\bar m,a),
                       \mathsf L(c\Rightarrow d,1)),
                  \widehat V_\chi^-(\bar m,v)),c\Rightarrow d,0\bigr).
\]
Thus the positive implication consumes evidence for the disjunction
\(\mathcal N^-\psi\lor\mathcal N^+\chi\), while its output
still has the original implication constructor.

For an unbounded existential with sign plus, let \(c\) be its original
matrix after substituting its other free parameters. Return
\[
     \mathsf{EI}(\widehat V_\psi^+(\bar m,a,v),i,c,\nu(a)).
\]
For an unbounded universal with sign minus, put
\(A=\mathsf{All}(i,c)\) and return
\[
 \mathsf{NI}\bigl(
  \mathsf B(\mathsf{FE}(\mathsf L(A,1),\nu(a)),
                     \widehat V_\psi^-(\bar m,a,v)),A,0\bigr).
\]
These are the only unbounded-quantifier cases on the stated domain:
the other two signs have an unbounded universal in their normal forms.

For the marked cases retain \(u,b,N,G,c,i,h,R_t,L\) from
\cref{def:p05-positive-evidence-builders}. Positive marked existential
uses
\[
 \mathsf{EI}\bigl(\mathsf{And}(L(\bar m,a),
                 \widehat V_\psi^+(\bar m,a,v)),
                                      i,G\mathbin{\&}c,\nu(a)\bigr).
\]
Negative marked universal, with \(A=\mathsf{All}(i,G\Rightarrow c)\),
uses
\[
 \mathsf{NI}\bigl(\mathsf B(
  \mathsf{MP}^{\rm row}(L(\bar m,a),
                    \mathsf{FE}(\mathsf L(A,1),\nu(a))),
                    \widehat V_\psi^-(\bar m,a,v)),A,0\bigr).
\]
For each of the other two cases form a separate list recursion
\[
\begin{aligned}
 \widehat Q_\theta^\epsilon(\bar m,z,0)&=0,\\
 \widehat Q_\theta^\epsilon(\bar m,z,Sk)&=
  \mathsf{cat}(\widehat Q_\theta^\epsilon(\bar m,z,k),
               [\widehat V_\psi^\epsilon(\bar m,k,\beta_*(z,k))]).
\end{aligned}
\]
Positive marked universal returns
\[
 \mathsf U_\diamond(c,i,h,u,b,R_t(\bar m),
                               \widehat Q_\theta^+(\bar m,z,N)).
\]
For negative marked existential put
\[
\begin{aligned}
 U&=\mathsf U_\diamond(\mathsf{Neg}(c),i,h,u,b,R_t(\bar m),
                               \widehat Q_\theta^-(\bar m,z,N)),\\
 E&=\mathsf{Ex}(i,G\mathbin{\&}c),\qquad
 y=1+\max(2,m(E)),\qquad
 K=\mathsf L(\mathsf{sub}(G\mathbin{\&}c,i,v_y),2).
\end{aligned}
\]
Return
\[
 \mathsf{NI}\left(\mathsf{EE}\left(\mathsf L(E,1),
  \mathsf B\left(\mathsf P_1(K),
      \mathsf{MP}^{\rm row}(\mathsf P_0(K),\mathsf{FE}(U,v_y))\right),
                                     i,G\mathbin{\&}c,y,1\right),E,0\right).
\]
The two displayed copies of \(K\) are separate leaves. Finally define
\[
 \widehat d_\theta^\epsilon(\bar m,z)\defeq
  I(e_\theta^\epsilon(\bar m,z),
                         \mathsf{pack}(\widehat V_\theta^\epsilon(\bar m,z)),0).
\]
All clauses insert only previously fixed signed-child descriptions.
\end{hbdefinition}

\begin{hbtheorem}{thm:p05-original-sigma-evidence}{proved=annex:p5.original-sigma-evidence}{Proof construction for every original fixed Sigma-one sentence}
For each fixed qualifying pair \((\theta,\epsilon)\), \(\PA\)
proves, with universal closure in \(\bar m,z\),
\[
\begin{aligned}
 \mathcal E_\theta^+(\bar m,z)&\to
  \mathsf{Safe}(\widehat V_\theta^+(\bar m,z),\theta^\nu(\bar m)),\\
 \mathcal E_\theta^-(\bar m,z)&\to
  \mathsf{Safe}(\widehat V_\theta^-(\bar m,z),
                                      \mathsf{Neg}(\theta^\nu(\bar m))),
\end{aligned}
\]
where only the line for that pair's sign is asserted. Each root is a
sentence. In particular, for every fixed closed \(\Sigma^0_1\) sentence
\(\sigma\) in its original syntax, \(\widehat d_\sigma^+\) is a
specified unary description and \(\PA\) proves
\[
 \forall z\,\bigl(\mathsf{Ev}_\sigma(z)\to
    \exists p\,(\Gamma_{\widehat d_\sigma^+}(z,p)\land
       \mathsf{Prf}^{19}_{\PA}(p,\overline{\ulcorner\sigma\urcorner}))\bigr).
\]
Consequently \(\PA\derives\sigma\to
\mathsf{Prv}^{19}_{\PA}(\overline{\ulcorner\sigma\urcorner})\).
The signed construction converts the evidence for the fixed normal form
to the original formula code at every formation step; its proof record
does not end at the normal-form code.
\end{hbtheorem}

\begin{hbdefinition}{def:p05-literal-introspection-instance}{The fixed original provability sentence and its builder}
Let \(P_{19}\) be the literal formula code of
\(\mathsf{Prv}^{19}_{\PA}(x)\) in
\cref{def:p05-selected-direct-pa-graph}, with its chosen free variable
index \(j\). Write its outer existential binder as \(i\); thus
\[
 P_{19}=\mathsf{Ex}(i,g_{19}),\qquad
 g_{19}=\ulcorner\Gamma_{d_{19}}(x_i,x_j,\bar1)\urcorner.
\]
These codes refer to the already chosen expanded graph, with its existing
bound indices. They do not call a fresh graph expansion after a numerical
input has been supplied.

For each fixed arithmetic sentence \(\varphi\), put
\[
\begin{aligned}
 n_\varphi&\defeq\ulcorner\varphi\urcorner,\\
 g_\varphi&\defeq\mathsf{sub}(g_{19},j,\nu(n_\varphi)),\\
 c_\varphi&\defeq\mathsf{sub}(P_{19},j,\nu(n_\varphi))
                         =\mathsf{Ex}(i,g_\varphi),\\
 \sigma_\varphi&\defeq
           \mathsf{Prv}^{19}_{\PA}(\overline{n_\varphi}).
\end{aligned}
\]
Thus \(c_\varphi=\ulcorner\sigma_\varphi\urcorner\). Let
\(\gamma_\varphi(x_i)\) be the formula with literal code
\(g_\varphi\). Its evidence and proof descriptions are the instances
of the preceding definitions for that fixed formula.

Define \(d^3_\varphi\) by the outer existential clause, explicitly
\[
 d^3_\varphi(z)\defeq
 I\left(e_{\sigma_\varphi}(z),
   \mathsf{pack}\left(
    \mathsf{EI}\bigl(
      V_{\gamma_\varphi}(\pi_0(z),\pi_1(z)),
                              i,g_\varphi,\nu(\pi_0(z))\bigr)\right),0\right).
\]
This is one separately fixed unary description for each
\(\varphi\). In particular its output root is \(c_\varphi\),
not \(n_\varphi\), a normalized substitute for \(c_\varphi\),
or a formula with newly chosen graph binders.
\end{hbdefinition}

\begin{hblemma}{thm:p05-literal-graph-evidence}{proved=annex:p5.literal-graph-evidence}{Every original graph and trace clause is covered}
Every \(\Gamma_d\) in the specified graph construction, and every
numeral instance of it, has a positive formation already in the fixed
marked negation normal form. This holds for the actual fully expanded
\(\Gamma_{d_{19}}\). Its evidence \(\mathcal E\) is therefore
the existing \(\mathsf{Ev}\), with the original fresh choices.

For each fixed arithmetic sentence \(\varphi\), \(\PA\) proves
\[
 \forall z\,\bigl(\mathsf{Ev}_{\sigma_\varphi}(z)\to
      \exists p\,(\Gamma_{d^3_\varphi}(z,p)\land
          \mathsf{Prf}^{19}_{\PA}(p,\overline{c_\varphi}))\bigr).
\]
The proof constructs the initial trace entry, every bounded step with its
two existential entry witnesses, the final entry and the outer trace
existential, using their original codes. A recursion counter zero and a
composition through a nullary function are included.
\end{hblemma}

\begin{hbtheorem}{thm:p05-direct-pa-d3}{proved=annex:p5.direct-pa-d3}{D3 for the literal selected PA predicate}
For every fixed arithmetic sentence \(\varphi\),
\[
 \PA\derives
  \mathsf{Prv}^{19}_{\PA}(\overline{\ulcorner\varphi\urcorner})
   \ \to\
  \mathsf{Prv}^{19}_{\PA}
       \left(\overline{\ulcorner
          \mathsf{Prv}^{19}_{\PA}
                    (\overline{\ulcorner\varphi\urcorner})
                         \urcorner}\right).
\]
Together with \cref{thm:p05-direct-pa-internal-mp}, this gives D1,
D2 and D3 for the same selected formula
\(\exists p\,\Gamma_{d_{19}}(p,x,\bar1)\).
No consistency assumption on \(\PA\) is needed for these derivability
conditions. The theorem is a scheme indexed by fixed sentences; it asserts
no total evaluation function for varying primitive recursive descriptions.
\end{hbtheorem}

\begin{hbexample}{def:p05-positive-evidence-zero}{A zero trace range with an original existential root}
Take a fixed formula
\(\tau=\forall x_i<0\,\exists x_k(x_k=x_k)\), with distinct
indices and with its original comparison binder retained. Evidence zero
satisfies its universal evidence clause, because the strict range is empty.
The constructor nevertheless builds the range-exhaustion implication to
bottom, the zero-stage case tree with target
\(\exists x_k(x_k=x_k)\), and the actual original-guard implication
and universal introductions. Its root is literally
\(\ulcorner\tau\urcorner\). No existential witness for that matrix
is selected in the empty fold. Replacing the bound by a weak zero bound
instead creates the one case at zero and requires its actual matrix
evidence. This distinguishes the two zero-bound constructions within the
same selected checker.
\end{hbexample}
